\documentclass[11pt, a4paper, answers]{exam}
\usepackage[a4paper, margin=2.5cm]{geometry}
\usepackage{fontspec}
% \usepackage[spanish, bidi=basic, provide=*]{babel}
% \babelfont{rm}{Noto Sans}
\usepackage{booktabs}
\usepackage{enumitem}

\title{\vspace{-2cm}\textbf{Evaluación de Geografía}}
\author{}
\date{}

\begin{document}
\maketitle
\vspace{-1.5cm}

\begin{center}
\fbox{\fbox{\parbox{0.9\textwidth}{\centering
Lee cuidadosamente cada instrucción y responde lo que se pide en cada sección.
}}}
\end{center}
\vspace{0.5cm}

\begin{questions}

\fullwidth{\section*{I. Tecnologías de la Información Geográfica}}

\question Clasifica las siguientes características indicando si pertenecen a un Sistema de Posicionamiento Global (\textbf{GPS}) o a un Sistema de Información Geográfica (\textbf{SIG}).
\begin{parts}
    \part \fillin[GPS][2cm] Permite posicionar con gran precisión puntos sobre la superficie de la Tierra.
    \part \fillin[SIG][2cm] Por medio de esta herramienta se puede utilizar la información registrada en bases de datos geográficos.
    \part \fillin[GPS][2cm] Sus siglas en inglés equivalen a Sistema de Posicionamiento Global.
    \part \fillin[SIG][2cm] Sus siglas significan Sistema de Información Geográfica.
    \part \fillin[GPS][2cm] Al inicio fue desarrollado con fines militares por Estados Unidos de América.
    \part \fillin[SIG][2cm] Son sistemas computarizados por medio de los cuales se pueden realizar mapas de diversos temas.
\end{parts}

\question Elige el sistema (GPS o SIG) que permite llevar a cabo las siguientes aplicaciones:
\begin{parts}
    \part \fillin[SIG][2cm] Monitorear especies animales en peligro de extinción.
    \part \fillin[GPS][2cm] Trasladarse de un lugar a otro en la vida cotidiana.
    \part \fillin[SIG][2cm] Considerar diferentes elementos para establecer un Área Natural Protegida.
    \part \fillin[SIG][2cm] Planear espacios urbanos con el objetivo de tener una mejor organización y funcionalidad.
    \part \fillin[SIG][2cm] Crear escenarios inexistentes para simular algún suceso, como una erupción volcánica o un huracán.
    \part \fillin[GPS][2cm] Trazar rutas de navegación terrestres, marítimas y aéreas.
    \part \fillin[GPS][2cm] Proporcionar datos correspondientes a la latitud, longitud y altitud (coordenadas).
\end{parts}

\question Completa el siguiente párrafo sobre las tecnologías utilizadas en el manejo de información geográfica con las palabras correspondientes:
\vspace{0.3cm}\\
En la actualidad existen diversas técnicas, métodos y \fillin[tecnologías][2.5cm] para obtener y \fillin[procesar][2cm] información del \fillin[espacio geográfico][3.5cm]. Dentro del tipo de tecnologías utilizadas para este fin se encuentran los equipos GPS (\fillin[Sistema de Posicionamiento Global][6.5cm]), los Sistemas de Información Geográfica (\fillin[SIG][1.5cm]), las imágenes de satélite y los vehículos aéreos tripulados de manera remota, también llamados \fillin[drones][2cm]. El GPS se utiliza para conocer la ubicación de los lugares, los objetos o las \fillin[personas][2.2cm]; por su parte, los SIG se componen de una computadora y un \fillin[software][2cm] o programa especializado útil para el análisis y procesamiento de la información obtenida de forma \fillin[satelital][2cm] o aérea.

\question Ordena del 1 al 7 el proceso utilizado por un GPS para el trazado de una ruta.
\begin{parts}
    \part \fillin[1][1cm] Tener un aparato receptor de señal GPS, por ejemplo, un celular.
    \part \fillin[2][1cm] Ubicar el lugar donde te encuentras.
    \part \fillin[3][1cm] Señalar el lugar donde deseas ir.
    \part \fillin[4][1cm] El celular por medio de un software especializado procesa la información.
    \part \fillin[5][1cm] Por medio de un algoritmo, el software calcula distancias y tiempos de recorrido.
    \part \fillin[6][1cm] El software señala la ruta a seguir en imagen.
    \part \fillin[7][1cm] El usuario se puede trasladar al lugar deseado.
\end{parts}

\newpage
\fullwidth{\section*{II. Dinámica de la Tierra: Vulcanismo, Relieve y Erosión}}

\question Indica si son Verdaderas (V) o Falsas (F) las afirmaciones con relación a la distribución de las regiones volcánicas.
\begin{parts}
    \part \fillin[F][1cm] El centro de México, específicamente la región fisiográfica del Sistema Volcánico Transversal, es una de las más activas del mundo.
    \part \fillin[V][1cm] Uno de los beneficios del vulcanismo son las aguas termales como atractivo turístico.
    \part \fillin[V][1cm] En la región del Atlántico se localiza una dorsal oceánica que atraviesa la Tierra de sur a norte y es donde se localizan los volcanes de Islandia.
    \part \fillin[F][1cm] A partir del vulcanismo se puede generar energía termoeléctrica.
    \part \fillin[V][1cm] México y Japón son países que aprovechan el vulcanismo para la generación de energía eléctrica.
\end{parts}

\question Completa los siguientes enunciados con la forma de relieve o concepto correspondiente:
\begin{parts}
    \part En las formas generales \fillin[del relieve][2.5cm] intervienen procesos internos y externos de la Tierra; unos intervienen en su formación y otros en su modelación.
    \part Las \fillin[montañas][2.5cm] pueden alcanzar alturas entre los 1 000 a 8 000 metros sobre el nivel medio del mar en promedio.
    \part Las \fillin[mesetas][2.5cm] o altiplanicies son áreas llanas o de terreno con poca pendiente situadas a cierta altitud.
    \part Las \fillin[llanuras][2.5cm] o planicies son zonas extensas de poca altitud.
    \part Las \fillin[depresiones][2.5cm] son partes del terreno hundidas que regularmente están ocupadas por agua.
\end{parts}

\question La erosión es un proceso externo de la Tierra que modifica el relieve. Elige la clase de erosión (\textbf{Eólica, Fluvial, Marítima}) a la que se refieren las descripciones.
\begin{parts}
    \part \fillin[Eólica][2cm] Ocurre principalmente en las regiones desérticas, como en el norte de México.
    \part \fillin[Fluvial][2cm] Por esta clase de erosión se modeló el Cañón del Sumidero en Chiapas.
    \part \fillin[Marítima][2cm] Tiene la capacidad de formar acantilados.
    \part \fillin[Eólica][2cm] Es la principal causa de que se formen las dunas y se modifiquen constantemente.
    \part \fillin[Fluvial][2cm] Las gotas de agua caen con gran velocidad y fuerza, lo que provoca el desgaste del suelo.
    \part \fillin[Marítima][2cm] Por causa de este tipo de erosión se ha formado el relieve característico de los Cabos en la Península de Baja California, en México.
    \part \fillin[Eólica][2cm] En las montañas, principalmente, no existe vegetación y sus laderas se ven modeladas por este tipo de erosión.
    \part \fillin[Fluvial][2cm] Cuando se presenta una lluvia torrencial y el agua acarrea el material en corrientes similares a ríos, se habla de una erosión...
    \part \fillin[Fluvial][2cm] Los fiordos, como los de Noruega, también son causa de la erosión...
    \part \fillin[Fluvial][2cm] Dicha clase de erosión también puede formar meandros en los ríos.
\end{parts}

\fullwidth{\section*{III. El Suelo y las Aguas Continentales}}

\question Completa los enunciados sobre el suelo:
\begin{parts}
    \part El suelo se forma debido a los procesos \fillin[externos][2cm] de la Tierra.
    \part Los suelos aportan servicios \fillin[ecosistémicos][3cm] que permiten la vida en la Tierra.
    \part El suelo es una \fillin[mezcla][2cm] de minerales, materia orgánica, agua, seres vivos y aire.
    \part El proceso de formación del suelo es muy \fillin[lento][2cm] y depende de varios factores químicos, \fillin[biológicos][2.5cm] y físicos.
\end{parts}

\question Indica si las afirmaciones sobre el daño o beneficio del suelo son Verdaderas (V) o Falsas (F).
\begin{parts}
    \part \fillin[F][1cm] El uso de agroquímicos en la agricultura beneficia los suelos.
    \part \fillin[V][1cm] El problema más común que afecta al suelo es la erosión.
    \part \fillin[V][1cm] La deforestación tiene un impacto negativo sobre el suelo.
    \part \fillin[F][1cm] La composta desfavorece la fertilización del suelo.
\end{parts}

\question Coloca las palabras para completar las afirmaciones sobre la clasificación de las aguas continentales.
\begin{parts}
    \part Se localizan en oquedades de la superficie terrestre, aunque también pueden existir \fillin[lagos][2cm] de tipo volcánico.
    \part Las aguas \fillin[subterráneas][3cm] se encuentran bajo la superficie terrestre.
    \part Los \fillin[glaciares][2.5cm] se localizan principalmente en regiones \fillin[polares][2cm], aunque en México se presentan en las cimas de algunos \fillin[volcanes][2.5cm].
    \part El llamado \fillin[cauce][2cm] principal de un río es alimentado por afluentes que integran las \fillin[cuencas][2.5cm] hídricas de México.
    \part Los mantos freáticos tienen su origen en la superficie terrestre por \fillin[filtración][2.5cm] del agua de lluvia.
\end{parts}

\fullwidth{\section*{IV. Categorías y Características del Espacio Geográfico}}

\question Elige la categoría de análisis del espacio geográfico (\textbf{Lugar, Paisaje, Región, Territorio}) que utilizarías para realizar los siguientes estudios.
\begin{parts}
    \part Jardín Escultórico de Edward James en Xilitla. \fillin[Lugar][3cm]
    \part Campos de viñedos a los pies de la Cordillera de los Andes en Chile. \fillin[Paisaje][3cm]
    \part Aprovechamiento de los recursos naturales para el turismo en Cayman Lodge Amazonia. \fillin[Lugar][3cm]
    \part Contaminación de las playas en la Costa Grande de Guerrero. \fillin[Región][3cm]
    \part Municipios y disponibilidad de agua en Sonora. \fillin[Territorio][3cm]
    \part Predominio industrial en Naucalpan, Estado de México. \fillin[Paisaje][3cm]
    \part Países que hablan alguna lengua de origen latino en América. \fillin[Región][3cm]
    \part Organización política de algunos países en Europa. \fillin[Territorio][3cm]
\end{parts}

\question Elige la palabra correcta del menú (lugar, territorio, paisaje, región) que complete los siguientes enunciados:
\begin{parts}
    \part Cayman Lodge Amazonia es un \fillin[lugar][2cm] en la Amazonia. Se localiza en el \fillin[territorio][2cm] de Perú.
    \part En las fotografías del \fillin[paisaje][2cm] de la selva amazónica, uno de los elementos principales es el río.
    \part La Amazonia abarca parte del \fillin[territorio][2cm] de Brasil, Bolivia y Colombia, entre otros países.
    \part Parte del proyecto de Cayman Lodge Amazonia es la protección de esta \fillin[región][2cm], que es considerada ``el pulmón principal'' de nuestro planeta.
    \part En el \fillin[paisaje][2cm] selvático de esta \fillin[región][2cm] natural predomina el componente natural de espacio geográfico.
    \part En este \fillin[lugar][2cm] llamado Cayman Lodge se puede convivir directamente con la naturaleza.
\end{parts}

\question Elige la característica del espacio geográfico (\textbf{Relación, Localización, Diversidad, Cambio, Distribución}) que le corresponde a cada descripción.
\begin{parts}
    \part Las coordenadas de Los Ángeles, California, en Estados Unidos de América, son: 34° 03' Norte; 118° 15' Oeste. \fillin[Localización][3cm]
    \part Santa Fe es una colonia de la Ciudad de México que hasta los años treinta era una zona de minas de arena, y en la actualidad es uno de los centros financieros más importantes del país. \fillin[Cambio][3cm]
    \part Las selvas, desiertos, bosques y mares forman parte de los ecosistemas que podemos encontrar en México. \fillin[Diversidad][3cm]
    \part En la Vega de Metztitlán, Hidalgo, la riqueza del suelo es favorecida por las crecientes del Río Venados, lo que permite una gran producción agrícola. \fillin[Relación][3cm]
    \part En los estados del norte del país se concentra el mayor número de maquiladoras. Dichos estados son: Baja California, Chihuahua, Tamaulipas, Coahuila de Zaragoza, Sonora y Nuevo León. \fillin[Distribución][3cm]
\end{parts}

\fullwidth{\section*{V. Representaciones Cartográficas y Símbolos}}

\question Relaciona el significado de los glifos representados en el mapa histórico de Mixquiahuala con la letra que corresponda según la figura provista por el profesor.
\begin{parts}
    \part \fillin[A][1.5cm] Línea de límite
    \part \fillin[B][1.5cm] Montaña
    \part \fillin[C][1.5cm] Iglesia
    \part \fillin[D][1.5cm] Gobernante
    \part \fillin[E][1.5cm] Vegetación
    \part \fillin[F][1.5cm] Río
    \part \fillin[G][1.5cm] Casas
\end{parts}

\question Coloca las palabras para completar las afirmaciones sobre fotografías aéreas.
\begin{parts}
    \part La \fillin[aviación][2cm] y la fotografía son inventos que permiten la toma de fotografías aéreas.
    \part En la actualidad, con las fotografías aéreas se pueden \fillin[reconocer][2.5cm] las zonas urbanas, los poblados, la vegetación, el trazo de las \fillin[calles][2cm] y todos los \fillin[componentes][2.5cm] del espacio geográfico.
    \part En la actualidad las fotografías aéreas están siendo sustituidas por fotografías tomadas con \fillin[drones][2cm] o por imágenes tomadas desde un \fillin[satélite][2.5cm].
\end{parts}

\question Observa la fotografía aérea agrícola y elige la respuesta correcta:
\begin{parts}
    \part ¿Qué representan los recuadros verdes en la fotografía aérea?
    \begin{oneparchoices}
        \choice Las localidades \CorrectChoice Los cultivos \choice La vegetación \choice Los cuerpos de agua
    \end{oneparchoices}
    \part ¿Qué representan los recuadros blancos?
    \begin{oneparchoices}
        \choice Las carpas \choice Las fábricas \choice Las casas \CorrectChoice Los invernaderos
    \end{oneparchoices}
    \part ¿Qué representan las líneas rectas?
    \begin{oneparchoices}
        \choice Los canales de riego \choice Las vías del tren \choice Los ríos \CorrectChoice Las carreteras
    \end{oneparchoices}
\end{parts}

\question Elige la categoría que corresponde a la característica o ejemplo de los tipos de símbolos (\textbf{Puntuales, Lineales, De área}).
\begin{parts}
    \part En esta categoría entran las superautopistas y se utilizan diferentes colores para indicar su importancia. \fillin[Lineales][2.5cm]
    \part En los mapas topográficos se utilizan para la representación de ciudades. \fillin[De área][2.5cm]
    \part Son ideales para la representación de regiones naturales, tipos de climas o cuerpos de agua como los océanos. \fillin[De área][2.5cm]
    \part Representan elementos aislados del espacio geográfico de importancia para la orientación de los lugares. \fillin[Puntuales][2.5cm]
    \part Representan límites de tipo natural o social, como las fronteras entre los países. \fillin[Lineales][2.5cm]
    \part Los hospitales utilizan este tipo de símbolos para su ubicación exacta en el mapa y poder ser ubicados rápidamente. \fillin[Puntuales][2.5cm]
\end{parts}

\fullwidth{\section*{VI. Análisis de Mapas y SIG}}

\question Indica si las afirmaciones sobre los usos o aplicaciones de los mapas son Verdaderas (V) o Falsas (F).
\begin{parts}
    \part \fillin[V][1cm] Se utilizan para el trazado de rutas y localización de lugares de desastre.
    \part \fillin[V][1cm] Los utilizan los gobiernos para delimitar predios y cobrar por el uso de suelo.
    \part \fillin[F][1cm] Con ellos se pueden ubicar de manera precisa y exacta las coordenadas de un punto sobre la superficie terrestre. (Nota: Se utilizan otras tecnologías de precisión como GPS para exactitud total).
    \part \fillin[V][1cm] Determinan rutas de evacuación en caso de una erupción volcánica.
    \part \fillin[F][1cm] Representan información en tiempo real de un desastre ocasionado por un sismo.
    \part \fillin[V][1cm] Gracias a ellos, se pueden determinar alturas aproximadas de elementos del relieve.
\end{parts}

\question De acuerdo con el mapa de Atracción y expulsión de migrantes, responde lo siguiente:
\begin{parts}
    \part ¿Qué nivel de atracción tienen los países de América del Norte?
    \begin{oneparchoices}
        \CorrectChoice Fuerte \choice Medio \choice Débil \choice Muy débil
    \end{oneparchoices}
    \part ¿En qué continente predominan los países con desplazamientos internos de migrantes?
    \begin{oneparchoices}
        \choice América \choice Europa \choice Asia \CorrectChoice África
    \end{oneparchoices}
    \part ¿Cuál es el principal destino de los brasileños?
    \begin{oneparchoices}
        \choice Australia \CorrectChoice Japón \choice Inglaterra \choice México
    \end{oneparchoices}
    \part Es uno de los principales destinos de los emigrantes de China.
    \begin{oneparchoices}
        \choice Francia \choice México \CorrectChoice Canadá \choice Reino Unido
    \end{oneparchoices}
    \part China presenta una expulsión de migrantes de tipo:
    \begin{oneparchoices}
        \CorrectChoice Fuerte \choice Medio \choice Débil \choice Muy débil
    \end{oneparchoices}
\end{parts}

\question Elige si las funciones de un SIG descritas son Verdaderas (V) o Falsas (F).
\begin{parts}
    \part \fillin[F][1cm] En Europa se dio a conocer el primer Sistema de Información Geográfica (SIG).
    \part \fillin[V][1cm] Los SIG integran un conjunto de herramientas para almacenar y analizar información acerca de las características del espacio geográfico.
    \part \fillin[F][1cm] Un mapa representa la realidad en la que vivimos.
    \part \fillin[F][1cm] Un SIG contiene toda la información geográfica de todo el mundo.
    \part \fillin[V][1cm] Los SIG tienen un funcionamiento que opera mediante diferentes módulos.
    \part \fillin[V][1cm] Los módulos de los SIG son de entrada, almacenamiento, salida y manejo de información.
    \part \fillin[F][1cm] La función de los SIG es la misma y se puede ver de igual forma en diferentes partes del mundo.
    \part \fillin[V][1cm] Los SIG guardan su información en una base de datos geográficos dentro de un sistema.
\end{parts}

\fullwidth{\section*{VII. Saberes Ancestrales y Medio Ambiente}}

\question Relaciona cada saber ancestral con el recurso que aprovecha y conserva al mismo tiempo.
\begin{parts}
    \part Agricultura de temporal: \fillin[Conocimiento del clima][5cm]
    \part Construcción de chinampas: \fillin[Canales lacustres][5cm]
    \part Cría de abejas meliponas: \fillin[Fauna nativa][5cm]
    \part Medicina tradicional: \fillin[Plantas medicinales][5cm]
\end{parts}

\question Elige si los siguientes saberes o técnicas son de origen \textbf{Ancestral} o \textbf{Moderno}.
\begin{parts}
    \part Relación de cultivos en la milpa. \fillin[Ancestral][3cm]
    \part Hidroponia. \fillin[Moderno][3cm]
    \part Uso de sistemas de riego por bombeo. \fillin[Moderno][3cm]
    \part Consumo de alimentos locales y de temporada. \fillin[Ancestral][3cm]
\end{parts}

\end{questions}
\end{document}